\section{Proofs for the normative-systems comparison}
\label{sec:p11-annex-normative-systems-comparison}

\begin{annexitem}{annex:p11.nsc.logical-checks}{thm:p11-nsc-logical-checks}{The source's example and the consequence/refutation distinction}
\begin{proof}
Universal elimination at the closed term Albert gives
\(\mathrm{worm}(\mathrm{Albert})\to\mathrm{warble}(\mathrm{Albert})\).
Conditional elimination with the second premise gives \(\varphi\).
With \(\neg\varphi\), conjunction introduction gives
\(\varphi\mathbin{\&}\neg\varphi\). In every model of the universal
premise, the worm extension is included in the warbling extension. The
worm premise puts \(a\) in the former; negative warbling excludes it
from the latter. No model can satisfy all three.

Deleting negative warbling allows the displayed finite model: its domain
contains every listed element, inclusion holds, and 5 belongs to both
extensions. Deleting the universal sentence allows domain \(\{5\}\),
worm extension \(\{5\}\), warble extension empty, and \(a=5\).
Deleting the worm sentence allows that same domain and denotation with
both extensions empty. The remaining universal is vacuously true and
negative warbling true. Thus every asserted repair has a model.

For arbitrary classical sentences, failure of \(\Gamma\entails\varphi\)
means precisely that a model makes every premise true and \(\varphi\)
false, hence makes \(\neg\varphi\) true. This is a model of
\(\Gamma\cup\{\neg\varphi\}\). Conversely any such model refutes
the consequence. If \(\Gamma\) has no model, there is no countermodel
for any conclusion, proving semantic explosion.

For natural deduction, an existing derivation of \(\varphi\) from
\(\Gamma\), under an additional assumption \(\neg\varphi\), yields
a contradiction. Conversely a contradiction derived with the extra
assumption \(\neg\varphi\) yields \(\varphi\) by classical reductio,
discharging that assumption. If it was unused, weakening permits the
same discharge. The proof uses that actual classical rule; it is not
asserted for every nonclassical calculus. The general soundness/completeness
theorems are supplied by the existing P2 proof interfaces, not inferred
merely from this example.
\end{proof}
\end{annexitem}

\begin{annexitem}{annex:p11.nsc.bayes}{thm:p11-nsc-bayes}{Multiplication, positivity and a nonidentified converse}
\begin{proof}
The definition gives \(\Pr(A,B)=\Pr(A\mid B)\Pr(B)\) when
\(\Pr(B)>0\), and \(\Pr(B,A)=\Pr(B\mid A)\Pr(A)\) when
\(\Pr(A)>0\). Intersection symmetry equates the left sides. Division
by the positive \(\Pr(A)\) proves the formula.

For an explicit comparison with one fixed likelihood, put mass on four
outcomes \((A,B),(A^c,B),(A,B^c),(A^c,B^c)\). The first law assigns
\((.1,.1,.4,.4)\); the second assigns \((.2,.2,.1,.5)\).
Both are nonnegative laws of total mass one. Both have
\(\Pr(A\mid B)=1/2\). The first has \(\Pr(B\mid A)=.1/.5=.2\),
whereas the second has \(.2/.3=2/3\). Every conditioning probability
used is positive. Thus the likelihood alone does not identify the converse.
The equalities relate one law's quantities; they contain neither a time
index nor a premise specifying which new information an agent learns.
No general behavioral update instruction follows from those equalities alone.
\end{proof}
\end{annexitem}

\begin{annexitem}{annex:p11.nsc.menu-representation}{thm:p11-nsc-menu-representation}{A full finite-menu proof}
\begin{proof}
Because \(C(\{x\})\) is a nonempty subset of a singleton, it equals
that singleton, proving reflexivity of \(\succeq_C\). Nonemptiness
of every pair choice implies that at least one of \(x\succeq_C y\)
and \(y\succeq_C x\) holds, proving completeness.

If \(x\in C(B)\), contraction to every pair \(\{x,y\}\) within
\(B\) gives \(x\succeq_C y\) for every \(y\in B\).
Conversely suppose all those comparisons hold. Enumerate the finite menu
as \(x,y_1,\ldots,y_m\). Each pair chooses \(x\). Expansion first
combines two such pair menus and then inductively adds every remaining
pair, proving that the full menu chooses \(x\). A singleton was already
handled. This proves equality with the maximal correspondence without
assuming transitivity.

Assume now transitivity. Indifference, meaning weak comparison in both
directions, is an equivalence relation. Weak transitivity makes comparison
between equivalence classes well defined; completeness orders the classes
and a strict comparison excludes a reverse comparison. There are finitely
many classes. List them from worst to best and give each member its class's
integer rank. This real function \(u\) represents all weak comparisons
and ties. The proved maximal-correspondence identity gives the argmax
formula, including multiple tied maximizers.

Conversely, in a finite complete transitive relation every nonempty menu
has a maximal class, so its maximal correspondence is nonempty. An option
maximal in \(B\) remains maximal in any sub-menu containing it, which
proves contraction. An option maximal in both \(A\) and \(B\) is
weakly at least as good as every member of their union, proving expansion.
The pair maxima include \(x\) exactly when \(x\succeq y\), so
the revealed relation recovers the given one. Finiteness and nonempty
menus were used essentially in this construction; it proves no unqualified
real representation on an arbitrary ordered domain.
\end{proof}
\end{annexitem}

\begin{annexitem}{annex:p11.nsc.pump}{thm:p11-nsc-pump}{Fee accounting and a complete counterpolicy}
\begin{proof}
Start at outcome \(P\) and wealth \(w\). The first accepted trade
leaves outcome \(M\) and wealth \(w-\Delta_1\); the second leaves
outcome \(P\) and wealth \(w-\Delta_1-\Delta_2\). Positivity makes
the loss positive. Induction on the number of accepted rounds gives
wealth \(w-k(\Delta_1+\Delta_2)\), provided every required round
remains affordable and accepted. No conclusion about unlimited feasible
repetition or an arbitrary initial wealth is smuggled into that premise.

For the counterpolicy, choose each singleton's member. On the remaining
free menus of \(\{P,M,D\}\), choose \(M\) from \(\{P,M\}\),
\(P\) from \(\{P,D\}\), \(M\) from \(\{M,D\}\), and
\(P\) from \(\{P,M,D\}\). This specifies every nonempty free
menu and violates contraction at the displayed pair/full menu. For any
offered transaction with a strictly positive fee, retain the current
outcome and wealth. This specifies refusal at every wealth level and
charges zero in both stages. The policy is a mathematical counterexample
to deriving willingness to pay from the free choices; it is not claimed
to represent a complete weak ordering or optimal human behavior.
\end{proof}
\end{annexitem}

\begin{annexitem}{annex:p11.nsc.worked}{thm:p11-nsc-worked}{All calculations and distinctions in the toy case}
\begin{proof}
The four masses sum to \((9+1+9+81)/100=1\).
Summing the first two gives \(\Pr(H)=10/100\); summing the first
and third gives \(\Pr(E)=18/100\). Thus
\(\Pr(E\mid H)=9/10\), \(\Pr(E\mid H^c)=9/90=1/10\),
and \(\Pr(H\mid E)=9/18=1/2\). All denominators are positive.
Each guess yields reward one in half the conditional outcomes and zero
in the other half, hence expectation \(1/2\). A perfect test always
classifies correctly and has net utility \(1-c\), so subtracting the
guess expectation gives \(1/2-c\). Its sign proves the three cases;
the two guesses tie throughout. At .2, test value is .8; at .6 it is .4.

The probability model is compatible with the logical conditional. In
\(H\)-worlds one can use a one-object model with Albert a warbling
worm; in \(H^c\)-worlds one can use Albert a nonworm who does not
warble, with the universal worm implication vacuously true. The report
event can occur in either kind of world. Therefore it is not an event
asserting that the actual premises hold. The deductive consequence was
proved for true premises, not for every occurrence of the fallible report.
No outcome or mental-process measurement enters these constructions.
\end{proof}
\end{annexitem}

\begin{annexitem}{annex:p11.nsc.sat-complete}{thm:p11-nsc-sat-complete}{Explicit verifier, local simulation circuit and clause reduction}
\begin{proof}
SAT belongs to NP: a certificate assigns one truth bit to every distinct
variable occurring in the encoded formula; evaluating the formula under
that assignment takes polynomial time in its length. Clauses of size at
most three admit the same check.

The verifier characterization is the same class as nondeterministic
polynomial-time acceptance. A fixed nondeterministic machine has only
finitely many possible transitions at a step; an accepting run of polynomial
length can be specified by a polynomial-length string of transition choices.
A deterministic verifier follows those choices, rejects invalid choices,
and checks the accepting state within a polynomial simulation bound.
Conversely a nondeterministic machine guesses the polynomially bounded
certificate and runs its deterministic verifier. Padding or a length
delimiter handles shorter certificates without changing either polynomial
bound. These two constructions justify the class identification used here.

Fix any verifier language and its deterministic verifier \(V\). We use
the standard one-tape verifier model with finite alphabet and state set,
one head and local transitions. Take polynomially bounded witness length
\(p(n)\). If the definition allows shorter witnesses, encode their
length in a prefix of polynomial size and pad the remaining bits; a verifier
checks the prefix and runs the original on that length's prefix. This
gives an equivalent fixed-length witness format. Increase the polynomial
time bound \(T\) so it also exceeds the initial tape length \(L\).
Arrange accepting/rejecting halts to remain in absorbing states thereafter.
The witness format and these changes remain polynomial.

The initial tape contains the fixed input \(x\), a separator, the witness
bits and blanks, with the head at the first cell. During \(T\) transitions
it cannot travel farther than \(T\) cells from its start. Retain all cells
from \(-T-1\) through \(L+T+1\), supplying blank neighbors outside
that interval. It contains the initial nonblank region and a margin beyond
every reachable head position.

Decorate each tape symbol by either no head or the current head state.
There is a fixed finite set of decorated symbols. Give each a fixed-width
binary code of \(b\) bits. For a legal configuration, the decorated value
of a cell at the next time is determined by the previous values of that cell
and its two neighbors: the transition writes at the old head and moves
the state decoration to one adjacent cell or retains it. This determines
a finite local function \(\mathsf F\) on legal triples. Give arbitrary
outputs to illegal triples/codes to make it total. For each output bit,
its finite truth table gives a Boolean expression: disjoin the conjunctions
of input literals for rows yielding one. Implement each disjunction through
De Morgan's identity using AND and NOT. Because \(V\) is fixed, \(b\)
and this local circuit's size are constants independent of \(x\).

Build a circuit layer for each of the \(T\) times and each retained cell,
feeding its three previous neighbor codes to this local circuit. At time
zero the input/blank/head codes are constants, and witness-symbol codes
are functions of the free witness bits. Boundary codes use the same blank
constants. Induction on time proves that on every witness the circuit
contains exactly the real computation: the initial configuration is legal,
and the local transition agrees with the next legal configuration.
Arbitrary values on illegal codes are never reached. At the final layer,
an OR of tests for an accepting-state code is the output. Absorption ensures
that accepting before \(T\) and accepting at \(T\) are equivalent.
The number of cells is \(L+2T+3\); thus circuit size and construction
time are polynomial, bounded by a constant times the number of layers
times retained cells. This is an actual finite construction, not a proposed
simulation that waits without a bound.

Convert the circuit to clauses. Introduce a Boolean variable \(z\) for
each gate-output wire; call its input-wire variables \(a_g,b_g\).
For an AND gate use the three clauses
\[
 (\neg z\lor a_g),\quad(\neg z\lor b_g),\quad
 (z\lor\neg a_g\lor\neg b_g).
\]
These force \(z=1\) exactly when both inputs are one: when \(z=1\)
the first two force both inputs, and when both inputs are one the third
forces \(z=1\). For a NOT gate use
\((z\lor a_g)\) and \((\neg z\lor\neg a_g)\), which force opposite
values. Constant-wire variables get their respective unit clauses.
Add the unit clause asserting that the final output is one. Each clause
has at most three literals, and at most three clauses per gate are added.
The resulting formula is polynomial in circuit size.

If a witness is accepted, assigning every wire its actual evaluated value
satisfies these clauses and the output clause. Conversely any satisfying
assignment obeys every gate equivalence; induction through the acyclic
circuit forces the output to be its evaluated value on that assignment's
witness bits. The output clause therefore yields an accepted witness.
The reduction maps \(x\) to one formula that is satisfiable exactly when
\(x\) belongs to the arbitrary original NP language. This proves NP-hardness
of the restricted clause problem and hence of SAT. Together with membership
it proves both completeness claims. If exactly three literals per clause
are required with repetitions allowed, repeat literals in shorter clauses
without changing satisfaction. No P-versus-NP separation is assumed.
\end{proof}
\end{annexitem}
